% GENERATED FILE. Edit canonical lessons, then run npm run generate:books.
% canonical-source-hash: fnv1a64:04d3791cc43e66e5
% canonical-lessons: IT-C423-crociera, IT-C423-cabina, IT-C423-marinaio, IT-C423-prua, IT-C423-poppa

\chapter{Sul traghetto}
\label{ch:it-sul-traghetto}
\hlchaptermodality{423}

\begin{chapteropening}
\textbf{By the end of this chapter:} \emph{I can talk about a voyage: the cruise, the cabin, the sailor, the bow and the stern.}

Complete the last lesson of chapter 423: I can talk about a voyage: the cruise, the cabin, the sailor, the bow and the stern.
\end{chapteropening}

\section[la crociera]{la crociera — the cruise}

\label{lesson:IT-C423-crociera}



\begin{quote}
Before the new one: say the Italian for to row, then the Italian for the canoe.
\end{quote}

\subsection*{You'll want to know: la crociera}
\textbf{la crociera} — \textquotedblleft{}the cruise\textquotedblright{}.

\begin{quote}
I miei genitori vanno in crociera.
\end{quote}

\begin{quote}
\emph{My parents are going on a cruise.}
\end{quote}

\emph{La velocità di crociera} — cruising speed.

\begin{cousinweb}
\emph{crociera} is built on \emph{croce} 'cross', which is inherited from Latin \emph{crux} (Lewis \& Short): first a ship crossing back and forth over a stretch of sea.
\end{cousinweb}

\begin{grammarlens}[title={What you've built}]
One more word for telling, explaining and describing at B1.
\end{grammarlens}

\subsection*{Your turn}
\begin{itemize}
\raggedright
  \item \emph{Say it:} \emph{la crociera}
  \item \emph{Say it:} \emph{la crociera}, once more
  \item \emph{Say it:} a sentence of your own with \emph{la crociera}
  \item \emph{Recall:} say the Italian for to row, then the Italian for the canoe, then say \emph{la crociera} again
\end{itemize}

\begin{tcolorbox}[breakable,colback=teal!4,colframe=teal!35!black,title={Before you move on}]
What does la crociera mean? (\textquotedblleft{}The cruise\textquotedblright{}.) Say it once more.
\end{tcolorbox}
\section[la cabina]{la cabina — the cabin; the booth}

\label{lesson:IT-C423-cabina}



\begin{quote}
Before the new one: say the Italian for the canoe, then the Italian for the cruise.
\end{quote}

\subsection*{You'll want to know: la cabina}
\textbf{la cabina} — \textquotedblleft{}the cabin; the booth\textquotedblright{}.

\begin{quote}
La nostra cabina ha una piccola finestra.
\end{quote}

\begin{quote}
\emph{Our cabin has a small window.}
\end{quote}

Also a booth: \emph{la cabina telefonica}.

\begin{cousinweb}
\emph{cabina} was borrowed from French \emph{cabine} (EtymDB); where the French word comes from is less certain.
\end{cousinweb}

\begin{grammarlens}[title={What you've built}]
One more word for telling, explaining and describing at B1.
\end{grammarlens}

\subsection*{Your turn}
\begin{itemize}
\raggedright
  \item \emph{Say it:} \emph{la cabina}
  \item \emph{Say it:} \emph{la cabina}, once more
  \item \emph{Say it:} a sentence of your own with \emph{la cabina}
  \item \emph{Recall:} say the Italian for the canoe, then the Italian for the cruise, then say \emph{la cabina} again
\end{itemize}

\begin{tcolorbox}[breakable,colback=teal!4,colframe=teal!35!black,title={Before you move on}]
What does la cabina mean? (\textquotedblleft{}The cabin; the booth\textquotedblright{}.) Say it once more.
\end{tcolorbox}
\section[il marinaio]{il marinaio — the sailor}

\label{lesson:IT-C423-marinaio}



\begin{quote}
Before the new one: say the Italian for the cruise, then the Italian for the cabin.
\end{quote}

\subsection*{You'll want to know: il marinaio}
\textbf{il marinaio} — \textquotedblleft{}the sailor\textquotedblright{}.

\begin{quote}
Un marinaio ci aiuta con le valigie.
\end{quote}

\begin{quote}
\emph{A sailor helps us with the suitcases.}
\end{quote}

From \emph{marina}, the sea and the navy.

\begin{cousinweb}
\emph{marinaio} is built on \emph{marina}, from \emph{mare} 'sea', which is inherited from Latin \emph{mare} (EtymDB).
\end{cousinweb}

\begin{grammarlens}[title={What you've built}]
One more word for telling, explaining and describing at B1.
\end{grammarlens}

\subsection*{Your turn}
\begin{itemize}
\raggedright
  \item \emph{Say it:} \emph{il marinaio}
  \item \emph{Say it:} \emph{il marinaio}, once more
  \item \emph{Say it:} a sentence of your own with \emph{il marinaio}
  \item \emph{Recall:} say the Italian for the cruise, then the Italian for the cabin, then say \emph{il marinaio} again
\end{itemize}

\begin{tcolorbox}[breakable,colback=teal!4,colframe=teal!35!black,title={Before you move on}]
What does il marinaio mean? (\textquotedblleft{}The sailor\textquotedblright{}.) Say it once more.
\end{tcolorbox}
\section[la prua]{la prua — the bow (of a ship), the prow}

\label{lesson:IT-C423-prua}



\begin{quote}
Before the new one: say the Italian for the cabin, then the Italian for the sailor.
\end{quote}

\subsection*{You'll want to know: la prua}
\textbf{la prua} — \textquotedblleft{}the bow (of a ship), the prow\textquotedblright{}.

\begin{quote}
Dalla prua si vede il porto.
\end{quote}

\begin{quote}
\emph{The harbour is in sight from the bow.}
\end{quote}

\emph{A prua} — at the front of the ship.

\begin{cousinweb}
\emph{prua} is inherited from Latin \emph{prōra} 'prow' (EtymDB), from Greek \emph{prōira}, probably through a dialect of the coast.
\end{cousinweb}

\begin{grammarlens}[title={What you've built}]
One more word for telling, explaining and describing at B1.
\end{grammarlens}

\subsection*{Your turn}
\begin{itemize}
\raggedright
  \item \emph{Say it:} \emph{la prua}
  \item \emph{Say it:} \emph{la prua}, once more
  \item \emph{Say it:} a sentence of your own with \emph{la prua}
  \item \emph{Recall:} say the Italian for the cabin, then the Italian for the sailor, then say \emph{la prua} again
\end{itemize}

\begin{tcolorbox}[breakable,colback=teal!4,colframe=teal!35!black,title={Before you move on}]
What does la prua mean? (\textquotedblleft{}The bow (of a ship), the prow\textquotedblright{}.) Say it once more.
\end{tcolorbox}
\section[la poppa]{la poppa — the stern (of a ship)}

\label{lesson:IT-C423-poppa}



\begin{quote}
Before the new one: say the Italian for the sailor, then the Italian for the bow.
\end{quote}

\subsection*{You'll want to know: la poppa}
\textbf{la poppa} — \textquotedblleft{}the stern (of a ship)\textquotedblright{}.

\begin{quote}
Ci sediamo a poppa, c'è meno vento.
\end{quote}

\begin{quote}
\emph{We sit at the stern, there is less wind.}
\end{quote}

\emph{Avere il vento in poppa} — to have the wind behind you; to do well.

\begin{cousinweb}
\emph{poppa} is inherited from Latin \emph{puppis} 'stern' (EtymDB).
\end{cousinweb}

\begin{grammarlens}[title={What you've built}]
One more word for telling, explaining and describing at B1.
\end{grammarlens}

\subsection*{Your turn}
\begin{itemize}
\raggedright
  \item \emph{Say it:} \emph{la poppa}
  \item \emph{Say it:} \emph{la poppa}, once more
  \item \emph{Say it:} a sentence of your own with \emph{la poppa}
  \item \emph{Recall:} say the Italian for the sailor, then the Italian for the bow, then say \emph{la poppa} again
\end{itemize}

\begin{tcolorbox}[breakable,colback=teal!4,colframe=teal!35!black,title={Before you move on}]
What does la poppa mean? (\textquotedblleft{}The stern (of a ship)\textquotedblright{}.) Say it once more.
\end{tcolorbox}
